% from ega0-14-fr.tex
\oldpage[0\textsubscript{IV-1}]{8}
car il en résulte que $\dim(U)=0$ pour un voisinage ouvert $U$ de $x$, et
comme $U$ est un espace de Kolmogoroff noethérien, $U$ est fini et discret.
\end{proof}

\begin{proposition}[14.1.11]
\label{0.14.1.11-fr}
La fonction $x\mapsto\dim_x(X)$ est semi-continue supérieurement dans $X$.
\end{proposition}

\begin{proof}
Il est clair que cette fonction est semi-continue supérieurement en tout
point où elle vaut $+\infty$. Supposons donc $\dim_x(X)=n<+\infty$~; alors,
la formule (14.1.4.1) montre qu'il existe un voisinage ouvert $U_0$ de $x$ tel
que $\dim(U)=n$ pour tout voisinage ouvert $U\subset U_0$ de $x$. Cela étant,
pour tout $y\in U_0$ et tout voisinage ouvert $V\subset U_0$ de $y$, on a
$\dim(V)\leq\dim(U_0)=n$ (14.1.4)~; on déduit donc de (14.1.4.1) que
$\dim_y(X)\leq n$.
\end{proof}

\begin{remark}[14.1.12]
\label{0.14.1.12-fr}
Si $X$, $Y$ sont deux espaces topologiques, et $f:X\to Y$ une application
continue, on peut avoir $\dim(f(X))>\dim(X)$~; on en a un exemple en prenant
pour $X$ un espace discret à 2 éléments $a$, $b$, pour $Y$ l'ensemble
$\{a,b\}$ muni de la topologie pour laquelle les ensembles fermés sont
$\varnothing$, $\{a\}$ et $\{a,b\}$~; si $f:X\to Y$ est l'application
identique, on a $\dim(Y)=1$ et $\dim(X)=0$. On notera que $Y$ est le spectre
d'un anneau de valuation discrète $A$, dont $a$ est l'unique point fermé et
$b$ le point générique~; si $K$ et $k$ sont le corps des fractions et le
corps résiduel de $A$, $X$ est le spectre de l'anneau $k\times K$ et $f$
l'application continue correspondant à l'homomorphisme
$(\varphi,\psi):A\to k\times K$, où $\varphi:A\to k$ et $\psi:A\to K$ sont
les homomorphismes canoniques (cf. IV, 5.4.3).
\end{remark}

\subsection{Codimension d'une partie fermée}
\label{subsection:0.14.2-fr}

\begin{definition}[14.2.1]
\label{0.14.2.1-fr}
Étant donnée une partie fermée irréductible $Y$ d'un espace topologique $X$,
on appelle \emph{codimension combinatoire} (ou simplement \emph{codimension})
de $Y$ dans $X$, et on note $\operatorname{codim}(Y,X)$ la borne supérieure
des longueurs des chaînes de parties fermées irréductibles de $X$, dont $Y$
est le plus petit élément. Si $Y$ est une partie fermée quelconque de $X$, on
appelle \emph{codimension de $Y$ dans $X$}, et on note encore
$\operatorname{codim}(Y,X)$ la borne inférieure des codimensions dans $X$ des
composantes irréductibles de $Y$. On dit que $X$ est
\emph{équicodimensionnel} si toutes les parties fermées irréductibles minimales
de $X$ ont même codimension dans $X$.
\end{definition}

Il résulte de cette définition que
$\operatorname{codim}(\varnothing,X)=+\infty$, la borne inférieure de la
partie vide de $\overline{\mathbb R}$ étant $+\infty$. Si $Y$ est fermé dans
$X$ et si $(X_\alpha)$ (resp. $(Y_\beta)$) est la famille des composantes
irréductibles de $X$ (resp. $Y$), tout $Y_\beta$ est contenu dans un
$X_\alpha$, et plus généralement toute chaîne de parties fermées
irréductibles de $X$ dont $Y_\beta$ est le plus petit élément est formée de
parties d'un $X_\alpha$~; on a donc
\[
  \operatorname{codim}(Y,X)=
  \inf_\beta\left(\sup_\alpha
  \bigl(\operatorname{codim}(Y_\beta,X_\alpha)\bigr)\right)
  \tag{14.2.1.1}\label{0.14.2.1.1-fr}
\]
où pour chaque $\beta$, $\alpha$ parcourt l'ensemble des indices tels que
$Y_\beta\subset X_\alpha$.

\begin{proposition}[14.2.2]
\label{0.14.2.2-fr}
Soit $X$ un espace topologique.
\begin{enumerate}
  \item[\rm(i)] Si $\Phi$ est l'ensemble des parties fermées irréductibles de
  $X$, on a
  \[
    \dim(X)=\sup_{Y\in\Phi}\bigl(\operatorname{codim}(Y,X)\bigr).
    \tag{14.2.2.1}\label{0.14.2.2.1-fr}
  \]
  